\section*{S2. Run-1 methods: local correction, fold-matched variant and hard labels}
\subsection{Novelty-damped local OOF residual correction}
\label{sec:naive}
With OOF residuals $r_j = y_j - s^{\mathrm{oof}}(x_j)$, the corrected prediction for a query $q$ is
\begin{equation}
\hat y(q) = s(q) + \alpha\,\frac{\sum_{j\in\mathcal N_k(q)} w_{qj}\,r_j}{\sum_{j\in\mathcal N_k(q)} w_{qj} + \kappa},
\qquad w_{qj}=\exp\!\big(-(d(q,j)/h)^2\big),
\label{eq:naive}
\end{equation}
where $\mathcal N_k(q)$ are the $k$ nearest training compositions under a composition distance $d$ (Euclidean on standardised Magpie vectors, or L1 on element fractions).
The pseudo-count $\kappa$ provides the novelty damping: when all neighbours are far, $\sum w$ is small and the correction vanishes.
For classification the shift is applied on the logit scale with probability residuals.
The tuple $(d,k,h,\kappa,\alpha)$ is selected on the training rows only, by scoring the corrected OOF predictions (leave-self-out neighbour sets) with the task metric over a fixed grid: $k\in\{5,10,20,40\}$, $h$ a multiple $\{0.5,1,2,4,8,16\}$ of the median nearest-neighbour distance, $\kappa\in\{0,0.1,0.5,2\}$, and $\alpha\in\{0.25,0.5,0.75,1\}$ (regression) or $\{0.1,0.25,0.5,1,2,4\}$ (classification).
The uncorrected stack ($\alpha=0$) is always a candidate and wins ties, so the procedure may decline to correct.

\subsection{Fold-matched correction}
\label{sec:fm}
The deployed stack $s=F^{\mathrm{full}}$ was trained on the very neighbours $j$ whose OOF residuals Eq.~\eqref{eq:naive} borrows.
Writing $F^{(-m)}$ for the stack whose experts excluded inner fold $P_m$, the naive transfer decomposes as $r_j = [r_j + F^{(-m(j))}(q) - F^{\mathrm{full}}(q)] + [F^{\mathrm{full}}(q)-F^{(-m(j))}(q)]$.
The second bracket, the \emph{absorption term}, measures how far the refit has already moved toward the neighbour, and the naive correction counts it a second time; for interpolating experts and near-duplicate neighbours it can be as large as $r_j$ itself.
The fold-matched variant transfers only the first bracket, $T_{qj}=r_j + F^{(-m(j))}(q) - F^{\mathrm{full}}(q)$, in place of $r_j$ in Eq.~\eqref{eq:naive}.
Its selection uses deployment-matched training instances: for a training row $i\in P_m$ and a neighbour $j\in P_{m'}$ ($m'\neq m$) the transfer is computed from the pair-excluded refit $F^{(-m,-m')}$, which saw neither $i$ nor $j$; neighbours in the same inner fold use $r_j$ directly.
This requires $J(J-1)/2=10$ extra fits per expert and fold, with the stack weights held fixed.
The variant was proposed by a structured ideation pipeline (ResearchStudio IdeaSpark), whose coherence check verified the algebraic decomposition on synthetic instances; we pre-registered it as a secondary arm, not as the primary method.
Relative to that proposal we simplified the damping map to the same $(\alpha,\kappa)$ family as Eq.~\eqref{eq:naive}, used the same two distances instead of the element mover's distance, and did not add TabPFN or MODNet experts, for reasons of CPU cost.

\subsection{Hard labels for classification}
Most classification entries on the Matbench leaderboard report a ``ROC-AUC'' equal to their balanced accuracy, which is what \texttt{matbench} computes when hard labels rather than probabilities are recorded.
For a like-for-like comparison we therefore also report balanced accuracy at a threshold chosen inside each training fold (maximising balanced accuracy on the arm's training OOF predictions), alongside the probability ROC-AUC.

